% 10-declarations.tex — Clause 10: Declarations
\chapter{Declarations}
\label{sec:10-declarations}
\index{declaration}

\pnum
This clause specifies top-level declarations.  Local variable declarations
are specified in \S\,\ref{sec:stmt.var}.  The grammar productions are in
Annex~A (\gramref{function-declaration}, \gramref{struct-declaration},
\gramref{enum-declaration}).

\section{Declaration attributes}
\label{sec:decl.attributes}
\index{attribute}

\pnum
A declaration may be preceded by zero or more attributes.  An attribute
has the form \tok{[name]} or \tok{[name(args)]}, where \textit{name} is
a recognized attribute identifier from the closed set in
Table~\ref{tab:attributes}, and \textit{args} (if present) is a
comma-separated list of expressions.  Multiple attributes on a single
declaration are written one per line.

\begin{laincode}
[fast_math]
func dot(a f64[], b f64[], n int) f64 { ... }

[private]
func internal_helper() int { ... }

[fast_math]
[private]
func hot_path() f64 { ... }
\end{laincode}

\pnum
Recognized attribute names (closed set, version 1.1):

\begin{longtable}{@{}lll@{}}
\caption{Recognized attribute names}\label{tab:attributes}\\
\toprule
\textbf{Attribute} & \textbf{Applies to} & \textbf{Purpose} \\
\midrule
\endfirsthead
\toprule
\textbf{Attribute} & \textbf{Applies to} & \textbf{Purpose} \\
\midrule
\endhead
\bottomrule
\endlastfoot
\tok{[fast\_math]} & any function & Enable FMA fusion (P4 opt-out, local). \\
\tok{[private]}    & any declaration         & Module-private visibility. \\
\end{longtable}

\begin{constraint}
  \textbf{Q-017.}
  An attribute name not listed in Table~\ref{tab:attributes} is
  \illformed{}~(\diagcode{E103}).  An attribute that omits the
  identifier (e.g.\ \texttt{[ ]} or \texttt{[(arg)]}) is
  \illformed{}~(\diagcode{E102}).
\end{constraint}

\begin{note}
  Attributes are recognized at declaration position only.  At expression
  position, \tok{[} and \tok{]} continue to delimit array and slice
  subscripts and array literals; the parser disambiguates by context.
\end{note}

\section{Struct declarations}
\label{sec:decl.struct}
\index{struct declaration}

\pnum
A struct declaration introduces a named product type with one or more
named fields.  Fields are listed in declaration order; order is significant
for construction and layout.

\begin{laincode}
type Point {
    x i32
    y i32
}
\end{laincode}

\pnum
A field may be annotated with \lain{mov}, making the field (and therefore
the containing struct) a linear type:

\begin{laincode}
type File {
    mov handle *FILE
}
\end{laincode}

\begin{note}
  A struct with zero fields (a \textit{unit} struct, \lain{type Empty \{\}})
  is permitted: it is a zero-sized marker type, constructed as \lain{Empty()}.
\end{note}

\section{Enum and ADT declarations}
\label{sec:decl.enum}
\index{enum declaration}
\index{ADT declaration}

\pnum
An enum declaration introduces a type whose values are drawn from a fixed,
named set of variants.  If every variant has no associated fields, the type
is a simple enum (\S\,\ref{sec:types.enum}).  If at least one variant has
associated fields, the type is an algebraic data type (ADT)
(\S\,\ref{sec:types.adt}).

\begin{laincode}
type Direction {              // simple enum (one variant per line)
    North
    South
    East
    West
}

type Shape {                                  // ADT
    Circle     { radius i32 }
    Rectangle  { width i32, height i32 }
    Point
}
\end{laincode}

\section{Function declarations}
\label{sec:decl.func}
\index{function declaration}

\pnum
A function declaration introduces a \lain{func}: a pure, terminating,
side-effect-free callable.  Its grammar production is
\gramref{function-declaration}.

\pnum
The signature consists of: the keyword \lain{func}, a name, an optional
list of type parameters (\S\,\ref{sec:14-generics}), a parameter list, and a
return type annotation.

\subsection{Parameters}
\label{sec:decl.func.params}
\index{parameter}

\pnum
Each parameter has a name, a type, and an ownership mode.  The ownership
mode determines how the argument is passed:

\begin{longtable}{@{}lll@{}}
\toprule
\textbf{Syntax} & \textbf{Mode} & \textbf{ABI (struct T)} \\
\midrule
\endfirsthead
\toprule
\textbf{Syntax} & \textbf{Mode} & \textbf{ABI (struct T)} \\
\midrule
\endhead
\bottomrule
\endlastfoot
\lain{name T}        & shared (read-only) & \texttt{const T*} \\
\lain{var name T}    & mutable            & \texttt{T*}       \\
\lain{mov name T}    & owned (transfer)   & \texttt{T} (by value) \\
\end{longtable}

\begin{note}
  For primitive types (integer, float, bool), the ABI passes the value
  directly regardless of ownership mode, as for primitive C arguments.
\end{note}

\pnum
A parameter whose declared type is the built-in \lain{type} --- written
\lain{name type}, with no \lain{comptime} prefix --- is a compile-time type
parameter; it is bound to a concrete type at every call site, by inference
from a value argument or explicitly (\S\,\ref{sec:14-generics}).

\subsection{Return type annotation}
\label{sec:decl.func.return}
\index{return type}

\pnum
The return type is written after the parameter list.  It may carry a value
range constraint:

\begin{laincode}
func abs(x int) int >= 0 { ... }
\end{laincode}

\begin{constraint}
  All return paths of a \lain{func} shall produce a value whose type and
  value range satisfy the declared return type and constraint.
\end{constraint}

\subsection{Purity constraints}
\label{sec:decl.func.purity}
\index{func purity constraints}

\begin{constraint}
  A function's \lain{effects} row is an upper bound on what its body does, and its
  default is $\emptyset$.  An effect the body has and the row does not name is
  \illformed{}~(\diagcode{E011} where the bound is $\emptyset$,
  \diagcode{E130} where a non-empty row understates it).  The three
  constraints below are consequences of that one rule, not separate rules.
\end{constraint}

\begin{constraint}
  A function shall not contain an unbounded \lain{while} loop --- one whose
  measure is neither given by \lain{decreasing} nor inferable --- unless its row
  names \lain{diverge}~(\diagcode{E011}/\diagcode{E082}).
\end{constraint}

\begin{constraint}
  A function shall not call one whose row names an effect its own row does
  not~(\diagcode{E011}).  The callee's \emph{introducer} is irrelevant: what is
  refused is the effect, not the keyword that declared it.
\end{constraint}

\begin{constraint}
  A function shall not participate in mutual recursion unless its row names
  \lain{diverge}~(\diagcode{E011}); termination cannot be established for a cycle
  no measure covers.  Direct self-recursion \emph{with} an inferred or given
  measure is permitted.
\end{constraint}

\section{The entry point}
\label{sec:decl.proc}
\index{entry point}

\pnum
\textbf{\lain{proc} has been removed.}  It was an introducer that granted effects
instead of stating them, and every effect it was needed for is now named by an
\lain{effects} row on an ordinary \lain{func} --- including \lain{diverge}, which
admits an unbounded loop.  The keyword remains \emph{reserved}, and using it is
\diagcode{E100} with the replacement spelled out, so old source fails with an
explanation rather than with "expected declaration".  The same applies to the
function-pointer type \lain{*proc(...)}, replaced by \lain{*func(...)} with a row
(\S\,\ref{sec:fn.distinction}).

\pnum
The entry point of every Lain program shall be declared as:

\begin{laincode}
func main() int { ... }          // with `effects io` if it performs any
\end{laincode}

\begin{constraint}
  A complete \emph{program} shall define exactly one \lain{main} returning \lain{int} as
  its entry point.  An individual module (translation unit) compiled on its own
  --- e.g.\ a library --- need not define \lain{main}: the reference compiler
  accepts a \lain{main}-less source file, and a program missing an entry point
  is diagnosed when the translation outputs are linked.
\end{constraint}

\section{Extern declarations}
\label{sec:decl.extern}
\index{extern declaration}

\pnum
An extern declaration introduces a C function into the Lain name space
without providing a body.  See \S\,\ref{sec:17-interop} for full semantics.

\begin{laincode}
extern func printf(fmt *u8, ...) int effects io
extern func strlen(s *u8) usize effects
\end{laincode}

\begin{constraint}
  Variadic parameters (\lain{...}) shall only appear in \lain{extern}
  declarations.  User-defined variadic functions are \illformed{}.
\end{constraint}

\section{Comptime declarations}
\label{sec:decl.comptime}
\index{comptime declaration}

\begin{note}
  Top-level \lain{comptime \{ ... \}} blocks --- evaluated by the compiler
  before any runtime code is emitted, to define compile-time constants or
  trigger \lain{compileError} assertions --- are a \emph{planned} feature not
  yet accepted by the reference implementation (a top-level \lain{comptime}
  block is presently \diagcode{E100}).  The \lain{comptime} keyword is
  nonetheless reserved, and \lain{comptime if} (\S\,\ref{sec:stmt.comptime-if})
  is available.
\end{note}
